\section{Source Pipeline and Run Provenance}
\label{app:pipeline}

\subsection{Calibration, confirmation, and source interpretation}
Removing eight development--test article overlaps leaves 2,651 eligible development questions.
A hash selects 176 questions from 50 articles for the ranking pilot.
The initial remainder contains 2,475 questions from 679 articles.
Earlier source inspections identify 35 further articles for exclusion before their first validation outcomes.
Removing their 127 questions gives the 2,348-question confirmation track.
The original split and exclusion records remain in the release.
The human track retains 830 positive questions from the published test collection.
Both tracks had earlier inspected results before the current revision.
The new adapter and selected ranking remain fixed throughout confirmation.

Current source calibration covers 86 articles, including 85 with eligible development questions.
It combines 50 pilot articles with 39 previously identified source paths, with three overlaps.
Ranking selection uses only the 176-question pilot.
An independent feature audit changes query metadata and remaps source identifiers without changing selections.
Ranking accepts the question string alone.
The title inventory comes from the source corpus.
Every ranking retains every corpus unit exactly once.

A year maps to January 1 through December 31 on the Gregorian ordinal axis.
A month maps to its actual first and last dates.
A fully stated day gives a singleton.
Tenures keep separate uncertain start and end events.
The adapter uses explicit clauses, coordinated tenures, biography constructions, and source-supported succession.
Unmatched or rejected text remains in the shared fallback partition.
All source words survive that partition.
Independent checks find no overlapping units and no rejected pair endpoints.
The development and test pools contain 33,433 and 34,480 units, respectively.

The independent calibration audit guides corrections before the adapter freeze.
Its final set contains 114 claims, including five succession starts and three directed pairs.
The full development corpus contains seven succession starts and those same three pairs.
The test corpus contains two succession starts and no pairs.
No rule invents a start date for an undated predecessor.
The Warburg example uses automatic extraction; its question period and ranking are illustrative.

A separate frozen audit samples 60 claims from 678 development claims outside calibration.
Selection uses a fixed identifier hash.
An independent model-assisted reviewer inspects full source paragraphs without benchmark questions or answers.
The reviewer accepts 52 events, rejects five, and marks three uncertain.
Eleven accepted starts omit a termination stated elsewhere in their paragraphs.
Of 29 sampled explicit-end records, 27 have accepted event grounding.
Rejected records contain reference errors, truncated names, or merged predicates.
The three uncertain records concern date attachment or combined roles.
These judgments do not change the frozen experiment.
They distinguish literal traceability, event interpretation, and lifecycle completeness.
They are model-assisted judgments, not human annotations or extraction-recall estimates.

\subsection{Reader protocol}
\label{app:reader}
The current model is Qwen2.5-3B-Instruct, Q4\_K\_M, at revision \texttt{7dabda4d}.
The release records the full revision and weight checksum.
The runtime is llama.cpp b11435 at commit \texttt{43fe9c6}.
The model uses the Qwen Research License; the runtime uses MIT.
Generation weights are excluded from the archive.
Pinned download scripts restore the official artifacts.

Prompts contain the visible question, original excerpts, source titles, date bounds, and temporal status.
Modeled excerpts also carry source-derived subject and relation labels.
Those labels preserve role context for short succession clauses.
No extracted answer-value field is added to the prompt metadata.
Temperature is zero, seed is 1729, and the output limit is 96 tokens.
Prompt caching and JSON grammar enforcement are disabled.
A tokenizer check verifies that every request fits the configured context.
Resource-only amendments precede completed question-answer outputs and preserve the original freezes and logs.
The final runtime uses four threads, a 2,048-token context, and polling disabled.
Its batch and microbatch sizes are both 512.

The parser accepts a complete JSON object or a single JSON code fence.
It rejects narrative repair, invalid field types, invalid excerpt indices, and truncated outputs.
Failures receive empty answer lists and remain in every denominator.
Normalization applies Unicode NFKC, case folding, and word-token extraction.
It preserves articles and compares deduplicated complete answer sets.
Strict set precision, recall, F1, and exact set accuracy use those normalized sets.
Token overlap remains a secondary measure.

The primary sample contains 60 questions from 55 articles.
Its 600 method conditions reduce to 193 distinct prompts.
The predeclared diagnostic selects every human query with different possible and guaranteed rendered contexts, up to twenty-four questions.
All eleven eligible questions enter that diagnostic, yielding 110 conditions and 45 distinct prompts.
They come from nine articles, four shared with the primary sample.
The cohorts have no question overlap and no shared requests.
Selection reads no answer labels or generated answers.
Full request hashes, raw responses, timing, and memberships support independent verification.
% Generated only from complete frozen-scoring outputs.
\begin{table}[t]
\centering\footnotesize
\setlength{\tabcolsep}{3pt}
\begin{tabular}{@{}lrrrr@{}}
\toprule
 & \multicolumn{2}{c}{Primary} & \multicolumn{2}{c}{Diagnostic}\\
Method & EM & F1 & EM & F1\\
\midrule
Original BM25 & 18.33 & 18.33 & 9.09 & 9.09\\
Original BM25 + year & 15.00 & 15.00 & 0.00 & 0.00\\
Title-aware ranking & 21.67 & 21.67 & 0.00 & 0.00\\
\quad + earliest completion & 20.00 & 20.00 & 18.18 & 24.24\\
\quad + midpoint completion & 20.00 & 20.00 & 0.00 & 6.06\\
\quad + latest completion & 20.00 & 20.00 & 0.00 & 6.06\\
\quad + interval control & 20.00 & 20.00 & 0.00 & 6.06\\
\quad + possible support & 20.00 & 20.00 & 0.00 & 6.06\\
\quad + guaranteed support & 20.00 & 20.00 & 18.18 & 24.24\\
Title-aware + year & 21.67 & 21.67 & 0.00 & 0.00\\
\bottomrule
\end{tabular}
\caption{Reader exact match (EM) and strict answer-set F1 (\%). Primary: 60 questions and 600 conditions. Diagnostic: 11 questions and 110 conditions. Unique generation parse failures: 75/193 and 13/45, respectively. Failures remain in both denominators.}
\label{tab:reader}
\end{table}


\paragraph{Complete reader outcomes.}
Table~\ref{tab:reader} reports both cohorts separately.
The diagnostic guaranteed-minus-possible F1 difference is 18.18 points, with article-cluster interval $[0.00,40.00]$.
All 45 paired method contrasts remain in each cohort's saved summary.
The primary temporal contrasts change no answer sets, including all 38 pairs with valid responses on both sides.
The diagnostic possible--guaranteed contrast changes three answer sets among six pairs with two valid responses.
Three further changes involve a response-format failure.
Earliest--latest completion gives the same counts and changed questions.
This selected diagnostic measures sensitivity among contexts already known to differ.

Across 238 distinct generations, 88 responses fail the frozen parser.
These comprise 60 JSON-format failures, six invalid citation lists, and 22 length-limit responses.
Every failed response remains scored with an empty answer set.
The primary ranking comparison changes 16 answer sets against original BM25 with recency.
Four changes occur among 23 pairs with two valid responses; twelve involve a parsing failure.
Separate descriptive audits report this breakdown for every method pair without changing the frozen scores.

\subsection{Ranking portability and selection cost}
A secondary analysis applies the certificate to all six ranking variants considered during calibration.
Its protocol is recorded after confirmation retrieval but before its own mechanism analysis.
It reads no answer labels and does not replace the selected main ranking.
Each variant has its own modeled-context denominator.
The release also reports fixed main and union cohorts for direct comparisons.
\begin{table*}[t]
\centering\small
\begin{tabular}{lrrrr}
\toprule
Ranking & Template stable/modeled & Human stable/modeled & Template tests & Human tests \\
\midrule
Global BM25 & 413/530 & 108/124 & 2.91 & 2.55 \\
Global BM25 + year & 670/810 & 191/224 & 3.00 & 2.63 \\
Title BM25 & 363/465 & 98/110 & 2.94 & 2.60 \\
Title-aware & 313/417 & 79/90 & 2.92 & 2.70 \\
Title BM25 + year & 529/646 & 155/175 & 2.95 & 2.65 \\
Title-aware + year & 463/569 & 146/166 & 2.77 & 2.60 \\
\bottomrule
\end{tabular}
\caption{Secondary certificate analysis across all six prespecified rankings. Each stability denominator contains modeled evidence under that ranking. Tests count temporal predicate calls within that modeled cohort.}
\label{tab:ranking-portability}
\end{table*}

The lazy selector agrees with full-mask evaluation for every parsed query.
Temporal-test counts include possible and guaranteed predicate calls.
Fallback excerpts require no predicate call.
The timing microbenchmarks exclude extraction, compilation, relevance ranking, rendering, and answer generation.
Concurrent CPU workloads affect their absolute timings.
The main text therefore reports operation counts.

\subsection{Earlier reader and extraction calibration}
Earlier Qwen2.5-1.5B-Instruct runs remain available under their original frozen settings.
A 24-question pilot requires 52 generations for 192 conditions.
A nine-question diagnostic adds 22 generations and reuses six.
The cohorts overlap on two questions, giving 31 distinct questions and 74 actual generations.
All methods reach 31.94\% strict set F1 on the pilot.
The earlier diagnostic changes four answer sets between earliest and latest completion, without improving aggregate strict F1.

Two exploratory learned extractors each process six calibration paragraphs before stopping.
The first returns 41 records, all missing required source quotes.
A schema-constrained variant returns 46 records; 45 fail literal grounding.
The remaining record has an incorrect holder--date association.
Neither extractor supplies claims to the reported confirmation runs.
Partial manifests retain all completed outputs, settings, and stopping decisions.

An earlier deterministic reproduction correction sorts BM25 token accumulation and stops after five accepted units.
Every one of 20,192 archived method contexts remains unchanged.
The original runner and equivalence records are preserved.

\subsection{Complete annual controls}
\begin{table}[t]
\centering
\footnotesize
\setlength{\tabcolsep}{3pt}
\begin{tabular}{lrr}
\toprule
Source-label property & Count & Rate (95\% CI) \\
\midrule
Multiple answers & 6,498/25,752 & 25.2 [24.2, 26.2] \\
Old first retained & 1,456/3,980 & 36.6 [35.1, 38.0] \\
Any answer retained & 6,050/6,961 & 86.9 [86.0, 87.8] \\
\bottomrule
\end{tabular}
\caption{Answer-set properties on held-out source keys. The last two rows condition on first-answer changes and complete-set changes, respectively. Rates and intervals are percentages.}
\label{tab:sourceaudit}
\end{table}

\begin{table}[t]
\centering
\small
\setlength{\tabcolsep}{4pt}
\begin{tabular}{lrrr}
\toprule
Policy & H@1 $\uparrow$ & R@5 $\uparrow$ & S@5 $\downarrow$ \\
\midrule
BM25 & 64.64 & 80.84 & 45.89 \\
BM25 + recent ties & 84.04 & 90.10 & 40.50 \\
Age weighting & 99.81 & 99.98 & 24.07 \\
Date-first ranking & 99.81 & 99.98 & 0.00 \\
Original predicate & 99.37 & 99.77 & 0.77 \\
Graphiti kernel & 99.78 & 99.93 & 0.00 \\
Query latest-value & 99.89 & 86.37 & 0.00 \\
Query latest-batch & 99.89 & 99.98 & 0.00 \\
\bottomrule
\end{tabular}
\caption{Annual source-label retrieval on 25,752 held-out questions. All values are percentages. H@1: first-result accuracy; R@5: answer recall; S@5: target-label stale exposure.}
\label{tab:fullsource}
\end{table}

The annual split and complete-source labels match the archived experiment.
Latest-value selection retains 99.89\% first-result accuracy but achieves 86.37\% complete answer recall.
Latest-batch selection reaches 99.98\% recall.
The suppression difference is 13.63 points, with key-cluster interval $[13.06,14.21]$.
Graphiti's matched policy kernel and direct-first invalidation produce identical endpoints across 45,866 observations.
These retained results require no repeated model inference.

\subsection{Earlier natural paragraph controls}
\begin{table}[t]
\centering
\small
\setlength{\tabcolsep}{3.5pt}
\begin{tabular}{lrrrr}
\toprule
& \multicolumn{2}{c}{Human} & \multicolumn{2}{c}{Templates} \\
Policy & H@5 & R@5 & H@5 & R@5 \\
\midrule
BM25 & 45.66 & 45.04 & 32.49 & 31.75 \\
Latest year & 54.70 & 54.02 & 51.86 & 50.89 \\
Year envelope & 49.16 & 48.33 & 34.44 & 33.67 \\
Explicit range & 44.10 & 43.53 & 32.11 & 31.39 \\
Range support & 44.58 & 44.00 & 38.50 & 37.63 \\
\bottomrule
\end{tabular}
\caption{TimeQA paragraph retrieval percentages. H@5 measures annotated-passage coverage; R@5 measures annotated-span recall. The tracks contain 830 human questions and 2,613 template questions.}
\label{tab:natural}
\end{table}

These paragraph runs use the published test passages and positive annotations.
Their human-question results were inspected before the source-unit experiments.
Paragraph and source-unit runs have different retrieval units and denominators.
